\clearpage\fancyhead[L]{\small KOMET.}\hypertarget{native-342-lower}{}\label{native-342-lower}
\begin{center}\textbf{KOMET.}\par\vspace{4pt}\hrule\vspace{4pt}\fontsize{10}{12}\selectfont Göttingische gelehrte Anzeigen. Stück 151. Seite 1497\,..\,1499. 1811 Sept. 21.\par\vspace{4pt}\hrule\vspace{4pt}\end{center}

\par\indent Den großen Kometen dieses Jahres, dessen Wiedererscheinen wir im 130. Stück unserer gelehrten Anzeigen [S. 327 d. B.] im Voraus angekündigt hatten, haben unsere Astronomen hier zum ersten Mal am 22. August in der späten Abenddämmerung gesehen. Eigentliche Beobachtungen waren erst später möglich, da die Gebäude der Stadt die Sicht von der Sternwarte nach Norden zu stark einschränken. Gemessen wurden die Abstände des Kometen von Fixsternen. Diese Methode eignet sich sehr gut für diesen hellen Kometen und liefert eine hohe Genauigkeit, besonders wenn die Beobachtungen vom Stativ aus gemacht werden. Von den auf diese Weise angestellten Beobachtungen sind bisher (10. September) allerdings erst die folgenden von Professor \textsc{Gauss} vollständig reduziert.

\Needspace{9\baselineskip}
\begin{center}\fontsize{10}{12}\selectfont\setlength{\tabcolsep}{4pt}\renewcommand{\arraystretch}{1.15}\begin{tabular}{ll|c|c|c}
\multicolumn{2}{c|}{1811} & Mittlere Zeit & Rektaszension & Nördliche Deklination \\\hline
Sept. & 4 & $8^{\mathrm u}\,28^{\mathrm m}\,47^{\mathrm s}$ & $158^\circ\;25'\;24''$ & $39^\circ\;18'\;2''$ \\
 & 6 & $8\;\;48\;\;38$ & $160\;\;\;23\;\;\;16$ & $40\;\;\;14\;\;\;16$ \\
 & 7 & $8\;\;57\;\;\;6$ & $161\;\;\;26\;\;\;10$ & $40\;\;\;41\;\;\;54$ \\
\end{tabular}
\end{center}

\par\indent Professor \textsc{Gauss} hat anhand dieser Beobachtungen die vorläufigen parabolischen Elemente verbessert, die er an der genannten Stelle mitgeteilt hatte, und folgende Ergebnisse gefunden.

\Needspace{13\baselineskip}
\begin{center}\fontsize{10}{12}\selectfont\setlength{\tabcolsep}{4pt}\renewcommand{\arraystretch}{1.15}\begin{tabular}{l}
Durchgang durch das Perihel 1811 Sept. 12.\\
\quad $5^{\mathrm u}\,21^{\mathrm m}\,15^{\mathrm s}$ mittlere Zeit in Göttingen\\
\end{tabular}\end{center}
\begin{center}\fontsize{9}{11}\selectfont\setlength{\tabcolsep}{2pt}\renewcommand{\arraystretch}{1.15}\begin{tabular}{llr}
Länge des Perihels & .\;.\;.\;.\;.\;.\;.\;.\;.\;.\;.\;.\;.\;.\;.\;. & $75^\circ\,17'\,34''$\\
Kleinster Abstand von der Sonne & .\;.\;.\;.\;.\;.\;.\;.\;.\;.\;.\;. & $1.04006$\\
Länge des aufsteigenden Knotens & .\;.\;.\;.\;.\;.\;.\;.\;.\;.\;.\;. & $140^\circ\,24'\,13''$\\
Neigung der Bahn & .\;.\;.\;.\;.\;.\;.\;.\;.\;.\;.\;.\;.\;.\;.\;.\;. & $73\;7\;16$\\
\end{tabular}\end{center}
\noindent Bewegung \emph{rückläufig}.

\par\indent Bisher gibt es keinen Hinweis darauf, dass die Bahn elliptisch ist.

\par\indent Um den Lauf des Kometen in den nächsten Monaten besser überblicken zu können, hat Professor \textsc{Gauss} aus den obigen verbesserten Elementen eine kurze Ephemeride berechnet. Daraus geht hervor, dass der Komet der Erde am 15. Oktober am nächsten kommt, dann aber noch fast ein Viertel weiter von ihr entfernt ist als die Sonne. Falls er überhaupt nur durch\SourceLeaf{343} reflektiertes Licht sichtbar ist, erreicht er seine größte Helligkeit am 7. Oktober. Er ist dann etwa fünfzehnmal so hell wie am 11. April. Am 3. Oktober nähert sich der Komet dem Stern $\eta$ im Großen Bären und am 3. Dezember dem Stern $\alpha$ im Adler bis auf weniger als einen Grad.

\Needspace{15\baselineskip}\noindent Lauf des Kometen während der letzten Monate dieses Jahres

\begin{center}\fontsize{10}{12}\selectfont\setlength{\tabcolsep}{3pt}\renewcommand{\arraystretch}{1.15}\begin{tabular}{ll|c|c|c|c}
\multicolumn{2}{c|}{1811} & \shortstack{Rektaszension} & \shortstack{Nördliche\\Deklination} & \shortstack{Abstand\\von der Erde} & \shortstack{Helligkeit} \\\hline
Sept. & 12\ $5^{\mathrm u}$ & $166^\circ\;54'$ & $42^\circ\;58'$ & 1.614 & 0.355 \\
 & 22\ $6$ & $182\;\;\;2$ & $47\;\;\;17$ & 1.443 & 0.432 \\
Okt. & 2\ $7$ & $203\;\;\;24$ & $49\;\;\;34$ & 1.306 & 0.490 \\
 & 12\ $8$ & $228\;\;\;45$ & $47\;\;\;16$ & 1.233 & 0.492 \\
\end{tabular}
\end{center}

\Needspace{15\baselineskip}\noindent Lauf des Kometen während der letzten Monate dieses Jahres

\begin{center}\fontsize{10}{12}\selectfont\setlength{\tabcolsep}{3pt}\renewcommand{\arraystretch}{1.15}\begin{tabular}{ll|c|c|c|c}
\multicolumn{2}{c|}{1811} & \shortstack{Rektaszension} & \shortstack{Nördliche\\Deklination} & \shortstack{Abstand\\von der Erde} & \shortstack{Helligkeit} \\\hline
Okt. & 22\ $0^{\mathrm u}$ & $251^\circ\;23'$ & $39^\circ\;56'$ & 1.248 & 0.421 \\
Nov. & 1\ $10$ & $267\;\;\;58$ & $30\;\;\;32$ & 1.353 & 0.310 \\
 & 11\ $11$ & $279\;\;\;38$ & $22\;\;\;18$ & 1.531 & 0.208 \\
 & 21\ $12$ & $288\;\;\;13$ & $14\;\;\;55$ & 1.744 & 0.138 \\
Dez. & 1\ $13$ & $294\;\;\;54$ & $9\;\;\;42$ & 1.984 & 0.093 \\
 & 11\ $14$ & $300\;\;\;24$ & $5\;\;\;49$ & 2.232 & 0.064 \\
 & 21\ $15$ & $303\;\;\;22$ & $3\;\;\;0$ & 2.426 & 0.048 \\
\end{tabular}
\end{center}

{\rightskip=0pt plus 1em\par\indent Sonderbar ist die Form des \emph{Schweifs}. Er biegt sich um den Kometen herum, ohne mit ihm verbunden zu sein, und läuft in zwei gegeneinander geneigte Äste aus. Von einem eigentlichen \emph{Kern} ist übrigens im Kometen keine Spur zu sehen.\par}
